%% LyX 2.4.4 created this file.  For more info, see https://www.lyx.org/.
%% Do not edit unless you really know what you are doing.
\documentclass[english]{article}
\usepackage[T1]{fontenc}
\usepackage[latin9]{inputenc}
\usepackage{enumitem}
\usepackage{amsmath}
\usepackage{amsthm}
\usepackage{cancel}
\usepackage{geometry}
\geometry{verbose,tmargin=1in,bmargin=1in,lmargin=1in,rmargin=1in}

\makeatletter

%%%%%%%%%%%%%%%%%%%%%%%%%%%%%% LyX specific LaTeX commands.
%% Because html converters don't know tabularnewline
\providecommand{\tabularnewline}{\\}

%%%%%%%%%%%%%%%%%%%%%%%%%%%%%% Textclass specific LaTeX commands.
\numberwithin{equation}{section}
\numberwithin{figure}{section}
\newlength{\lyxlabelwidth}      % auxiliary length 

%%%%%%%%%%%%%%%%%%%%%%%%%%%%%% User specified LaTeX commands.
\usepackage{fancyhdr}
\usepackage{lastpage}
\pagestyle{fancy}
\usepackage{enumitem}
\usepackage{ifthen}
%sets math fonts to be more readable
\everymath=\expandafter{\the\everymath\displaystyle}
%\DeclareMathSizes{10}{10}{10}{10}
\usepackage{multicol}

\usepackage{tikz}
\usetikzlibrary{plotmarks}
\usepackage{pgfplots}
\pgfplotsset{compat=newest}

\usepackage{calc}
\usepackage[nomessages]{fp}

\usepackage{advdate}    % Advancing/saving dates
\usepackage{datetime}   % Dates formatting
\usepackage{datenumber} % Counters for dates
% Added by lyx2lyx
\setlength{\parskip}{\smallskipamount}
\setlength{\parindent}{0pt}

\makeatother

\usepackage{babel}
\begin{document}
\renewcommand{\author}{Dr.\,James Ashe}

\newcommand{\classNum}{132}

\newcommand{\classSec}{C and K}

\newcommand{\nameType}{Name:}

\newcommand{\examType}{Worksheet 6.1 and 6.2}

\renewcommand{\date}{\formatdate{7}{4}{2014}}

%%First page's header%%%%%%%%%%%%%%%%%%%%%%
\fancypagestyle{firststyle} %%header settings for first pagr
{
	\fancyhead[L]{MTH \classNum{}(\classSec{}) \author{}\\
	\textbf{\examType{}} \date{}}
	\fancyhead[C]{\textbf{\nameType}}
	\setlength{\headsep}{14pt}
}
%%%%%%%%%%%%%%%%%%%%%%%%%%%%%%%%%%%%%%%%%%

%%Next page's header%%%%%%%%%%%%%%%%%%%%%%%
\setlist{nolistsep}
\lhead{\classNum{}(MTH \classSec{})} 
\chead{\textbf{\nameType}} 
\cfoot{\thepage{}~of~\pageref{LastPage}} 
\setlength{\headsep}{5pt} 
%%%%%%%%%%%%%%%%%%%%%%%%%%%%%%%%%%%%%%%%%%%

%%Various settings; see the preamble for other settings%%%%%
%%\setenumerate[0]{leftmargin=12pt,itemindent=0pt,labelwidth=0pt}
\renewcommand{\labelenumi}{\textbf{\arabic{enumi}.}}
\renewcommand{\labelenumii}{\textbf{\alph{enumii}.}}
\thispagestyle{firststyle}
%%%%%%%%%%%%%%%%%%%%%%%%%%%%%%%%%%%%%%%%%%
\newcommand{\xmax}{10}
\newcommand{\xmin}{-10}
\newcommand{\xstep}{0} %must be xmin+n
\newcommand{\ymax}{10}
\newcommand{\ymin}{-10}
\newcommand{\ystep}{0} %must be ymin+n

\newcommand{\xspace}{\vspace{.75cm}}

\textbf{Decide whether each of the following is  statement. If it
is a statement, decide whether or not it is compound.}
\begin{enumerate}[resume]
\item Nashville is the capital of Tennessee.\xspace
\item Got milk?\xspace
\item China does not have a population of more than 1 billion, and it is
a country.\xspace
\item If I get an ``A'' then I will celebrate.\xspace
\end{enumerate}
\textbf{Write a negation of each statement.}
\begin{enumerate}[resume]
\item This is not the time to complain.\xspace
\item $y\leq5$\xspace
\item $p$\xspace
\item $p\wedge q$\xspace
\end{enumerate}
\textbf{Fill in the following truth tables}.\begin{multicols}{2}

\begin{tabular}{cc|c}
$p$ & $q$ & $p\wedge q$\tabularnewline
\hline 
\hline 
T & T & \tabularnewline
\hline 
T & F & \tabularnewline
\hline 
F & T & \tabularnewline
\hline 
F & F & \tabularnewline
\hline 
\end{tabular}\columnbreak

\begin{tabular}{cc|c}
$p$ & $q$ & $p\vee q$\tabularnewline
\hline 
\hline 
T & T & \tabularnewline
\hline 
T & F & \tabularnewline
\hline 
F & T & \tabularnewline
\hline 
F & F & \tabularnewline
\hline 
\end{tabular}\end{multicols}

\textbf{Let $p$ represent a false statement and $q$ represent a
true statement. Find the truth value of each compound statement.}
\begin{enumerate}[resume]
\item $\sim p$\vfill
\item $p\vee\sim q$\vfill
\item $\sim(p\vee\sim q)$\vfill\newpage
\end{enumerate}
\textbf{Let $p$ represent a true statement, and $q$ and $r$ represent
false statements. Find the truth value of each compound statement.}
\begin{enumerate}[resume]
\item $q\vee\sim r$\vfill
\item $(\sim p\wedge q)\vee\sim r$\vfill
\end{enumerate}
\textbf{Let $p$ be the statement ``$2>y$'', $q$ be statement
``$8\leq6$,'' and $r$ be the statement ``$19\leq19$.'' Find
the truth value of each compound statement.}
\begin{enumerate}[resume]
\item $\sim p\vee\sim r$\vfill
\item $(p\wedge q)\vee\sim r$\vfill
\item $(\sim r\wedge q)\vee\sim p$\vfill
\end{enumerate}

\end{document}
